% 斯特凡·巴拿赫（综述）
% license CCBYSA3
% type Wiki

本文根据 CC-BY-SA 协议转载翻译自维基百科\href{https://en.wikipedia.org/wiki/Stefan_Banach}{相关文章}

\begin{figure}[ht]
\centering
\includegraphics[width=6cm]{./figures/13afb1dae6b98ddb.png}
\caption{} \label{fig_STFbnh_1}
\end{figure}
斯特凡·巴拿赫（波兰语发音：[ˈstɛfan ˈbanax]；1892年3月30日—1945年8月31日）是一位波兰数学家，被普遍认为是20世纪最重要、最具影响力的数学家之一。他是现代泛函分析的奠基人之一，也是利沃夫数学学派的核心成员之一。他的主要著作是1932年出版的《线性算子论》（法语原名 Théorie des opérations linéaires），这本书是关于泛函分析一般理论的第一部专著。

巴拿赫出生于克拉科夫的一个具有戈拉尔（Góral，高地人）血统的家庭。他自幼对数学表现出浓厚兴趣，常在课间休息时解数学题。完成中学教育后，他结识了数学家雨果·施泰恩豪斯，二人于1919年共同创立了波兰数学学会，并后来创办了学术期刊《数学研究》。1920年，他在利沃夫理工学院获得助教职位，1922年升任教授，1924年成为波兰学术院院士。巴拿赫还是利沃夫数学学派的共同创立者之一，该学派汇聚了波兰战间期（1918–1939）最著名的一批数学家。

以巴拿赫命名的数学概念包括：巴拿赫空间、巴拿赫代数、巴拿赫测度、巴拿赫–塔斯基悖论、哈恩–巴拿赫定理、巴拿赫–施泰恩豪斯定理、巴拿赫–马祖尔博弈、巴拿赫–阿劳格鲁定理、巴拿赫–萨克斯性质以及巴拿赫不动点定理等。
\subsection{生平}
\subsubsection{早年生活}
斯特凡·巴拿赫于1892年3月30日出生在克拉科夫的圣拉撒路综合医院，当时该地属于奥匈帝国。他出身于一个戈拉尔族的罗马天主教家庭，随后由父亲为他施洗。巴拿赫的父母分别是斯特凡·格雷切克和卡塔日娜·巴拿赫，二人都来自波德哈莱地区。格雷切克是奥匈帝国军队驻克拉科夫的士兵，而关于巴拿赫的母亲，人们所知甚少。根据洗礼记录，她出生于博鲁夫纳，是一名家庭女佣。

不同寻常的是，斯特凡的姓氏取自母亲，而不是父亲，但他的名字则继承自父亲。由于军队规定禁止格雷切克这样的士兵（当时军衔为列兵）结婚，而母亲又贫困无力抚养孩子，这对男女决定让孩子由亲友抚养。斯特凡生命的最初几年由祖母照料，但当祖母生病后，格雷切克安排他由克拉科夫的弗朗齐丝卡·普沃娃和她的侄女玛丽亚·普哈尔斯卡抚养。年幼的斯特凡视弗朗齐丝卡为养母，而玛丽亚则如同姐姐一般。在早年，巴拿赫由法籍知识分子尤利乌什·米恩（Juliusz Mien）辅导，他是普沃娃家的朋友，曾移居波兰，以摄影与将波兰文学翻译成法语为生。米恩教授巴拿赫法语，并极有可能在他早期的数学启蒙中起到了鼓励作用。
